\documentclass[12pt,a4paper]{report}

% ====================
% PACKAGES & SETTINGS
% ====================
\usepackage{mathptmx} % Times New Roman Font
\usepackage[utf8]{inputenc}
\usepackage{setspace}
\usepackage{graphicx}
\usepackage{amsmath} 
\usepackage{float} 
\usepackage{url}
\onehalfspacing % 1.5 Line Spacing

\usepackage{titlesec}
\usepackage{geometry}
\geometry{a4paper, margin=1in}

% ====================
% COMPLIANCE FONT SIZES
% ====================
% Chapter 1 (Heading) - Font Size 16 - Bold
% Chapter Title - Font Size 14 - Bold
\titleformat{\chapter}[display]
  {\normalfont\fontsize{16}{19}\bfseries\centering}{\MakeUppercase{\chaptertitlename\ \thechapter}}{10pt}{\fontsize{14}{17}\bfseries}

% Sub headings - Font Size 12 - Bold
\titleformat{\section}
  {\normalfont\fontsize{12}{14}\bfseries}{\thesection}{1em}{}
\titleformat{\subsection}
  {\normalfont\fontsize{12}{14}\bfseries}{\thesubsection}{1em}{}

\begin{document}

% ====================
% FRONTMATTER
% ====================
\pagenumbering{roman}

\chapter*{SYNOPSIS}
\addcontentsline{toc}{chapter}{SYNOPSIS}

Traditional methods for evaluating software engineering have been based on static evaluations of code along with the use of unit testing paradigms. These methods typically do not allow for the assessment of a developer's ability to address issues such as dynamic board state changes; finite resources; or the potential for adversarial logic execution.

The coding game platform proposed by EMERGENT represents a paradigmatic shift in how engineers are evaluated. Unlike static puzzles, developers will be required to create bot algorithms using competing logic that operate in real time within a turn-based environment. This platform was designed to be built atop a "Room First" architecture so that it can isolate tenants locally, and is constructed of a React.js Single Page Application (SPA), with a language agnostic Monaco Editor. Thus, this platform allows participants to develop their ideas without having to install anything into their computers. However, since the platform executes untrusted remote code there is a significant risk involved in providing access to remote code. Therefore, the platform utilizes a queue to process compilations of logic and passes them off to a decoupled Referee Daemon. This daemon processes these matches and guarantees Zero Trust Isolation from the host machine through the utilization of throttled ephemeral Docker containers.

Unlike other grading systems, where developers receive a binary pass/fail grade, the EMERGENT platform utilizes a multidimensional statistical model called Glicko-2. This model continuously evaluates the engineering resilience of developers. Furthermore, the platform uses a decoupled Visualizer to render deterministic JSON artifacts created from log data.
\newpage
\setcounter{tocdepth}{1} % Limits ToC to Chapters and Sections (1.1, 1.2), hiding Subsections (1.1.1)
\tableofcontents
\newpage
\listoffigures
\newpage
\listoftables
% \newpage

\chapter*{LIST OF ABBREVIATIONS}
\addcontentsline{toc}{chapter}{LIST OF ABBREVIATIONS}

\vspace{0.1cm}
\noindent
\renewcommand{\arraystretch}{1.5}
\begin{tabular}{@{}ll}
\textbf{API} & Application Programming Interface \\
\textbf{AST} & Abstract Syntax Tree \\
\textbf{CORS} & Cross-Origin Resource Sharing \\
\textbf{DOM} & Document Object Model \\
\textbf{E2E} & End-to-End \\
\textbf{FPS} & Frames Per Second \\
\textbf{GUI} & Graphical User Interface \\
\textbf{HMR} & Hot-Module Replacement \\
\textbf{HTTP} & Hypertext Transfer Protocol \\
\textbf{IDE} & Integrated Development Environment \\
\textbf{IPC} & Inter-Process Communication \\
\textbf{JSON} & JavaScript Object Notation \\
\textbf{OOM} & Out Of Memory \\
\textbf{ORM} & Object-Relational Mapping \\
\textbf{OS} & Operating System \\
\textbf{RAM} & Random Access Memory \\
\textbf{RD} & Rating Deviation \\
\textbf{REST} & Representational State Transfer \\
\textbf{SPA} & Single Page Application \\
\textbf{UUID} & Universally Unique Identifier \\
\textbf{WSGI} & Web Server Gateway Interface \\
\end{tabular}
\newpage

\pagenumbering{arabic}

% ====================
% CHAPTER 1: INTRODUCTION
% ====================
\chapter{INTRODUCTION}

\section{PROJECT OVERVIEW}
EMERGENT is a Local Area Network (LAN), advanced, fully autonomous, single-server coding game platform, created to provide an extremely interactive environment for developers to develop their skills in developing algorithmic solutions. Unlike most traditional coding environments, EMERGENT creates a "sandbox" environment where the developer's algorithms or heuristics interactively compete with each other in real-time. A "turn-based" game type is developed where players' "bots" interact and attempt to defeat each other. An example of one such game would be "Tron". 

The architecture of EMERGENT was built with a "room-first" design philosophy. Each "room" has its own namespace. Therefore, the user does not have to contend with latency due to global servers and leader boards. By creating a room first, the friction involved with deploying new tournaments is greatly reduced. Additionally, it enables tournament administrators or coding club leaders to easily create their own private competitions utilizing consumer-grade hardware. The administrator can create a room, select a pre-defined set of rules for the game being played, and receive a cryptographic key ("PIN") along with a room identifier (UUID). When participants enter into the lobby they will need to supply a username and upon entering the lobby they may immediately start writing their competing algorithm. The development of each competitor's algorithm occurs within an integrated web-based Integrated Development Environment (IDE). 

Once a competitor submits their algorithm via the internet; the backend system's multi-threaded engine queue's the submitted scripts, isolates them from each other, compiles them, runs them competitively and returns rankings according to the results. The rankings are generated using the well-established Glicko-2 formula. The lobby leaderboard is updated in near real-time based on wins, losses and/or draws. The server also generates replays, displays console outputs and logs exceptions errors such as runtime errors caused by syntax failure or OOM (Out-of-Memory) kernel interrupt errors or exceeding time limits for execution.

\section{PROBLEM STATEMENT AND MOTIVATION}
Industry standards for measuring engineering competency and logical resiliency focus almost entirely on static local assessment techniques. The traditional teaching methods used today for DSA (Data Structures \& Algorithms) such as LeetCode, HackerRank, etc., teach students only how to pass / fail based on a statically defined set of inputs, that test the code solely against the correct mathematical solutions.

Static problem solving does allow students to learn basic Data Structures \& Algorithms (DSA). However, Static Problem Solving is unable to mimic the dynamic interaction of real world applications when dealing with adverse inputs, changing states, and limited hardware resources. Students are rarely educated on how to write graceful degradation into their code, manage state effectively over time, or respond to uncontrolled external influences.

This is where EMERGENT comes in. It addresses this exact issue through the creation of a Bot Battle arena that allows users' code to be tested not just for accuracy, but also for strategy and long term viability. A Bot that cannot respond appropriately to an unforeseen change in its Board State will fail catastrophically. This failure is directly attributed to the poor design choices made by the developer regarding exceptions.

Furthermore, through the use of a proprietary Rating Algorithm, EMERGENT provides users with empirical ratings relative to other localized developers allowing for a direct comparison of coding skill. Allowing for an environment that encourages iterative development of a user's code to achieve excellence in all aspects of coding.

\section{SCOPE AND OBJECTIVES}
The primary objective of the EMERGENT system is to architect a platform that completely abstracts away the complexities of local environment setups for the end-user, while providing an impenetrable execution sandbox for the host application. 

The explicit goals of the minimum viable product (MVP) encompass:
\begin{itemize}
    \item \textbf{Frictionless Onboarding:} Requiring exactly zero local installations (compilers, interpreters, or dependencies) for the student. Everything must operate seamlessly via a modern web browser.
    \item \textbf{Guaranteed Determinism:} Ensuring the Referee engine simulation executes the game loop linearly and identically every single time when given the same random seed and input code, yielding absolute algorithmic fairness.
    \item \textbf{Zero-Trust Code Execution:} Implementing an architecture where malicious code (e.g., recursive fork bombs, remote network pings, or endless infinite memory arrays) submitted by a user cannot structurally escape the bounded container or crash the host application.
    \item \textbf{Rapid Analytical Feedback:} Generating highly granular \texttt{GameState.json} logs post-match that instantly allow the user to visualize the match sequentially frame-by-frame on a Canvas UI, fundamentally resolving "Black Box" testing issues by showing exactly why and where an algorithm failed logically.
\end{itemize}

\section{SYSTEM CONFIGURATIONS AND PREREQUISITES}
Because the platform simulates an entire virtual computational fleet using background daemon workers alongside an active web server, the hardware recommendations are categorized by the operational load expected on the host device.

\subsection{Hardware Requirements}
For optimal and stable performance of the EMERGENT host tournament server, particularly prioritizing the asynchronous background execution engine, the following hardware constraints are evaluated:
\begin{itemize}
    \item \textbf{Processor (CPU):} Minimum multi-core processor (Intel Core i5, AMD Ryzen 5, or equivalent ARM chips). High thread concurrency is required because the Referee routinely spawns Python sub-processes heavily reliant on internal stream pipe routing.
    \item \textbf{Memory (RAM):} 8 GB minimum, though 16 GB is heavily recommended. Because the Docker daemon imposes a strict 512MB memory boundary per executing container, a queue processing multiple matches concurrently will dynamically reserve and flush massive logical memory blocks.
    \item \textbf{Storage Volumes:} 40 GB of free space Minimum, specifically allocated to support the dense, localized SQLite database string stores, the Docker Engine virtual image snapshots, and localized user bot source code repositories. Solid State Drives (NVMe/SSD) are strongly recommended over Hard Disk Drives (HDD) to reduce container boot I/O latency.
\end{itemize}

Client machines connecting to the web interface require merely a standard modern machine with an updated web browser and a stable LAN connection to the host server.

\subsection{Software Requirements}
\begin{itemize}
    \item \textbf{Operating System:} A Linux Kernel-based OS (Ubuntu 22.04 LTS or Debian 12 preferred) is highly recommended over Windows due to Docker's explicit reliance on native \texttt{cgroups} and namespaces for container isolation.
    \item \textbf{Language Runtimes:} Python 3.10+ (for backend orchestration and multi-processing) and Node.js v18+ (for executing Vite frontend bundle compilation scripts).
    \item \textbf{Containerization Platform:} Docker Engine v24+ configured to bind directly to the background worker interface for process sandboxing.
    \item \textbf{Client Access:} Google Chrome version 90+, Mozilla Firefox 89+, or any equivalent modern browser offering strict support for HTML5 Canvas rendering and WebSocket background polling.
\end{itemize}

\section{TECHNOLOGY OVERVIEW AND DEVELOPMENT STACK}
EMERGENT is built upon a modern, modular, and containerized technology stack optimized fundamentally for rapid prototyping and heavily isolated computational execution. 

\subsection{Backend Routing and Orchestration Layers}
The core of the API layer is comprised of \textit{\textbf{Flask}}, which is a lightweight Python Web Server Gateway Interface (WSGI). Flask primarily functions as a stateless gateway to route incoming HTTP requests to specific APIs. Flask creates Room session schema dynamically based on the structure of canonical HTTP request content. Lastly, it will function as an interface that serves API endpoints for the frontend's data processing. To remove any large-scale structural complexity associated with LAN hosted tournament operations, the data layer is using a highly normalized relational SQLite database. Using this relational database model eliminates the need for any of the common ORM frameworks. This enables the system to directly execute raw SQL commands locally. As such, the system has eliminated the normal latency associated with a database round trip over a network when making write commits at high frequency during matchmaking events.

\subsection{Execution Sandbox and The Referee}
The main API thread of the core competitive engine uses an asynchronous Redis queue backed by continuously running python worker threads. Each time a user submits a code string, that string is dynamically bound to a physical file and mounted as a volume inside a docker container. Since the Docker system serves as the final authority on isolating processes; we are able to enforce additional restrictions that isolate our server from potential malicious submissions through setting internal engine flags for example --network=none --cpus="1.0" along with specific memory allocations. Once isolated in the container, the referee will establish dynamic communication with the referee's simulation process using IPC stream (standard I/O pipe). As each abstract game is implemented using local python plugins within a registry of local plugins (for example /server/games/tron/plugin.py), the addition of new types of games will allow us to add them to the overall architecture of the umpire without having to make changes to the basic routing mechanisms.

\subsection{Frontend Presentation and The Developer Experience}
This portal was created using a highly optimized Single Page Application (SPA) that utilizes JavaScript to leverage the \textit{\textbf{React.js}} framework for building components. This SPA is then booted through a \textit{\textbf{Vite}} development pipeline to allow for millisecond level Hot Module Replacements (HMR).

In order to create a seamless coding experience, this application embeds the \textit{\textbf{Monaco Editor}}, which is the same open source declarative editor structure used in Microsoft's Visual Studio Code. When Monaco is imported into the browser, it allows the browser to provide language agnostic abstract syntax trees (AST) for code formatting, algorithmically folds code and provides syntax highlighting on a line-by-line basis for languages including Python, Java and C++.

Finally, custom built React HTML5 components will capture the raw JSON data from the \textit{GameState.json} output from the Docker backend and use this data to programmatically rebuild the 2-dimensional representation of the game state allowing users to interactively view each step of a match inside the native web portal.
\section{REPORT ORGANIZATION}
The remainder of this report is meticulously organized to sequentially document the architectural design, algorithmic implementation, and systemic testing of the EMERGENT platform.
\begin{itemize}
    \item \textbf{Chapter 2: System Analysis} explores the problem boundaries in detail, outlining standard limitations of existing static coding solutions, identifying explicit non-functional limits, and thoroughly cataloguing the requirements necessary.
    \item \textbf{Chapter 3: System Design} charts the formal architecture patterns using rigorous Unified Modeling Language (UML) specifications, modeling actor behaviors, asynchronous process workflows, and exact underlying entity-relationship (ER) mapping rules optimized for local Room tenancy.
    \item \textbf{Chapter 4: System Implementation} dissects the deeply integrated technical challenges regarding how exactly the Docker streams interact safely, the exact Mathematical representations of the Glicko-2 Elo system, and the deployment of the Web Monaco IDE.
    \item \textbf{Chapter 5: Testing} catalogues the extensive White-Box and Black-Box quality assurance methodologies enforced, defining specific unit verifications and Sandbox edge-case resilience tests mapping server exceptions.
    \item \textbf{Chapter 6: Conclusion and Future Scope} finalizes the report by reflecting heavily on the objectives successfully realized during the MVP, and maps heavily towards theoretical Cloud-Native expansions, including Kubernetes processing layers and Data Analytics integration.
\end{itemize}

% ====================

% ====================
% CHAPTER 2: SYSTEM ANALYSIS
% ====================
\chapter{SYSTEM ANALYSIS}

\section {EXISTING SYSTEM} 
As stated above, the current state of the art in automated programming assessment are mainly based on two categories of platforms: static algorithmic problem solver and game-based coding platform. An overview of the existing literature and commercial products demonstrate there exists a large gap in the area of testing of real world engineering resiliency.

\subsection {Static Algorithmic Problem Solvers}
Platfroms such as LeetCode and HackerRank have become the de facto industry standards for recruiting technical talent and measuring their abilities. These platforms follow a static analysis and unit testing methodology whereby the user's code is compared against a predetermined set of input/output scenarios. While very effective at demonstrating an individual's basic understanding of data structures and algorithms (DSA), platforms such as LeetCode/HackerRank do not simulate day-to-day engineering challenges. Each submission to either platform results in a binary outcome (pass/fail) and is also a fully isolated event; therefore the submitted code will never interact with any external systems or actively malicious entities. As a result, these platforms cannot measure an individual's ability to design systems that can adapt to changing circumstances or evaluate holistic design decisions.

\subsection {Game-Based Coding Platforms}
In order to encourage ongoing usage in static platforms, services such as CodinGame integrated gamification into their platforms through visual feedback and "bot battle" type of games to promote continued learning. Although previous research has demonstrated that these types of environments increase cognitive engagement with the material, these platforms suffer from what we refer to as the "relative skill" issue. That is most game-based platforms allow users to be ranked relative to one another without providing a reference point (i.e., a calibrated AI opponent) to measure a user's absolute skill. Therefore, it is often difficult for recruiters or educators to accurately assess an applicant's overall technical ability, i.e., a high ranking may simply indicate that the user had an advantage due to playing against weaker opponents.

\subsection {Gap Between Dynamic and Static Testing} 
There is a fundamental distinction between determining whether or not an application has been correctly written (static) versus determining whether or not a system is resilient (dynamic). Technical literature has emphasized that while static testing provides early defect detection capabilities it does not provide adequate capability to detect defects resulting from runtime anomalies associated with race conditions, resource depletion, or hostile input. In addition, since EMERGENT employs a dynamic execution engine that simulates resource utilization and evaluates how well an application performs when subjected to hostile input, it directly addresses this deficiency.

\section {PROPOSED SYSTEM }EMERGENT identifies and fills key shortcomings in the current landscape with a “Room first” model employing a dynamic multi-agent sandbox. The EMERGENT model fundamentally changes how evaluations are performed from discrete assessments against static baselines to continuous performance metrics generated solely from head-to-head competition. Thus the EMERGENT platform rewards code resiliency and long-term thinking (strategic foresight) rather than short-term memory.
\subsection {Deployment of a deterministic execution engine}The main technological contribution will be the development of a "referee system" container which is able to execute un-trusted user code securely. A "Warm Pool" architecture design is used to greatly minimize the time required for containers to start up and therefore provide an inherent scalability model for low-latency/high-security code execution. As the referee engine strictly orders the input to each Turn sequentially, every game is deterministic and can be reproduced using mathematical methods.
\subsection {A Standardized Skill Assessment Framework}In contrast to other systems where users receive either a pass/fail assessment or a single numerical score (e.g., a rating), the proposed system generates a Continuous Elo Rating (utilizing glicko-2) for each user. In addition to providing users/recruiters with a means of assessing the relative engineering maturity of a candidate, the proposed system also introduces "tiered benchmark bots" (calibrated Tier 1-3 ai opponent bots). Users/recruiters can use these bots to assess the absolute engineering maturity of candidates.
\subsection {Enhanced Debugging Capabilities (“the flight Recorder”) }The proposed platform sets a new standard for debugging in competitive environments by producing detailed Turn-by-Turn replay logs of matches in JSON format. These logs allow developers to perform white box testing on their algorithms, enabling them to view a match in detail — Turn by Turn — within the frontend visualizer. Developers may determine precisely at what point/why/when a particular strategy failed – something completely absent in standard Black Box testing platforms.

\section{REQUIREMENTS SPECIFICATION}

\subsection{Functional Requirements}
The functional requirements of the EMERGENT platform are categorized into four distinct modules:
\begin{itemize}
    \item \textbf{User Interface \& Bot Management:} The system provides an embedded, browser-based IDE (powered by Monaco Editor) featuring exact syntax highlighting and line numbering for Python, Java, and C++. It enables bot submission, versioning, and displays comprehensive personal statistics (Elo, win/loss ratios).
    \item \textbf{Matchmaking \& Orchestration:} Utilizing an asynchronous priority queue via Redis, incoming match requests are handled without interrupting main server threads. The matchmaker employs skill-based pairing to automatically match participants with similar ratings and manages real-time updates to the room-focused leaderboard.
    \item \textbf{The Execution Engine (Referee):} The core module exclusively processes matches inside strictly isolated Docker environments, protecting the host system. It maintains a "Warm Pool" of paused containers to reduce match latency below 500ms and guarantees 100\% deterministic simulation sequences.
    \item \textbf{Replay \& Visualization:} Following a match, the engine must render a precise \texttt{GameState} architectural log. The frontend React visualizer parses this artefact to construct a fully interactive GUI replay player including pause, rewind, and tick-inspection mechanics.
\end{itemize}

\subsection{Non-Functional Requirements}
\begin{itemize}
    \item \textbf{Security \& Isolation:} All user-submitted logic is fiercely sandboxed inside Docker containers without host network or sensitive file system access. Mandatory Resource Quotas forcefully terminate executions exceeding 512MB of RAM or 1.0 seconds of CPU duration per game turn to prevent system Denial of Service (DoS) attacks.
    \item \textbf{Performance \& Latency:} Utilizing the warm pool structure, match request-to-execution latency must remain under 2 seconds. Global workspace UI and editor initialization must load under 1.5 seconds on standard network configurations, and subsequent leaderboards must push updates within 5 seconds of match completion.
    \item \textbf{Reliability \& Determinism:} Reproducibility is non-negotiable. Re-evaluating the exact same bot code pairing under the exact same random game seed must always yield a bit-for-bit identical final simulation outcome. Furthermore, the orchestrator should strive for a 99\% uptime availability rate during active tournament sessions.
    \item \textbf{Scalability:} To support high-velocity tournament contexts, the system must smoothly orchestrate a minimum of 50 concurrent matches (equating to 100 active participant containers) concurrently via effective queue pooling mechanics.
    \item \textbf{Usability:} Runtime or Syntax execution exceptions must be accurately captured, parsed out of the unformatted standard error streams, and delivered to the user interface in a human-readable format indicating the specific logic point of failure.
\end{itemize}

% ====================
% CHAPTER 3: SYSTEM DESIGN
% ====================
\chapter{SYSTEM DESIGN}

The System Design stage is the most important part in developing the architecture, or overall plan, for a system to meet all the defined technical and business needs. In this stage, the development team had to be extremely thoughtful about their decisions on how to build the EMERGENT platform's architecture because it needed to overcome the very large operational difference (as well as the performance limitations) between general-purpose web applications and high-security/low-latency environments where a completely deterministic game-execution engine will be running autonomously. The rest of this chapter includes details of the complete model of the EMERGENT Platform, the design pattern used, the architecture of the API Layer, the sequence of asynchronous data processing during code evaluation, and the relational database that underlies a foundationally robust structure that strictly adheres to the "Room-First" multi-tenancy rules.

\section{USE CASE DIAGRAM}
A Use Case diagram provides a high-level visual representation of the interactions between the primary external actors and the EMERGENT platform. It outlines the scope of the system and the functional boundaries of each user role within the architectural boundary. 

In this system, the interactions are primarily divided between two actors:
\begin{itemize}
    \item \textbf{Participant (Developer):} The main user of the platform represents people that will actually be competing against each other in the algorithmic competitions (the main event). These users want to create, improve and ultimately implement algorithms into bots to perform as efficiently as possible during these tournaments. Participants go through the same basic process when interacting with the platform for developing their algorithms. First, they join one of the available sessions via a valid Room Code. This establishes their presence in the current matchmaking state. Next, the participant uses the built-in Monaco IDE to develop, test, debug and refine their algorithm while being contained in a safe sandboxed environment. At some point, once their development cycle is complete, they will then submit their compiled code to the systems execution queue for scheduling. Once scheduled, their code executes and they can view the "GameState.json" graphically interactive replay. The replay provides all the information about what occurred in real time and allows the participant to review their strategy and logically identify areas where they can make improvements to better compete at future events. Ultimately, it is up to the participants to continually develop, evaluate, debug and repeat until they have achieved competitive success.
    
    \item \textbf{System Umpire (Execution Engine):} The automated backend agent responsible for orchestration. It seamlessly pulls queued submission matches, provisions the isolated Docker sandboxes, enforces the rigid CPU/RAM constraints, calculates the post-match Glicko-2 participant ratings, and records the match outputs directly to the database.
\end{itemize}

\begin{figure}
    \centering
    % Placeholder for the Mermaid.js Use Case Diagram Image
    \includegraphics[width=13cm]{pictures/usecase_project_diagram.png}
    \caption{Use Case Diagram depicting actor interactions in the EMERGENT platform}
    \label{fig:use_case_diagram}
\end{figure}

\section{SYSTEM ARCHITECTURE PATTERN}
EMERGENT has a decoupled, asynchronous, hybrid micro-service architecture. Although a fully deployed service mesh was considered unnecessary overhead for the local scope of a tournament server, using a purely monolithic architecture will be a disaster due to the blocking nature of local compilation execution.

\subsection{Presentation Layer (Client)} 
The Presentation Layer is the only access point for humans and resides entirely within the client’s web-browser. It is built on top of the SPA (Single Page Application) model. This layer takes care of managing local-state (for example; active editor text buffers and cached replay of matches) while at the same time does not require complete page refreshes from the server. In addition, component modularity allows each UI fragment (Visualizer, Leader Board, Monaco IDE) to operate independently with the exception of when their respective data payloads change.

\subsection{Orchestration & API Layer (Backend)}
The API Layer is the central point-of-interaction and exists solely to act as a router. Due to its design being intentionally “thin”, the API Layer does not perform game logic but instead functions almost exclusively as a gatekeeper and message broker. Initially, it validates requests, enforces syntactical correctness of room codes and parses incoming multipart/form-data (users’ zipped bot code). Once a user’s code is successfully written to disk, the API Layer immediately passes the computationally intensive task off to the next tier by sending a “match job” through either an internal pub/sub mechanism or polling to the background worker queue(s), allowing the web-server to maintain very low-latency responses for end-users.

\subsection{Asynchronous Execution Layer (Background Worker Pool)}
Completely isolated from all API endpoints, the Execution Layer provides a pool of daemon process running in the background. The Execution Layer utilizes a Worker Queue Design Pattern. Workers are running in an endless synchronous event-loop where they pull match-jobs from the data-store. Since moving sandbox instantiation into this layer protects the web-server from hanging if a student writes an infinite loop, this layer explicitly interfaces with the docker-daemon as a low-level service to rapidly spin-up/tear-down ephemeral containers.

\subsection{Data Persistence Layer} 
Finally, the bottom-most layer is the Data Persistence Layer which consists of both SQLite relational database and file-system. Because of the high-read/high-write frequency of the match-maker (which continuously updates dynamic Glicko-2 ratings) — SQLite was used over higher-weighted cloud databases (i.e., PostgresSQL) to remove any potential round-trip network latency. In terms of file structure, there is a significant focus placed upon sandboxing user-directories, so that no Python bot in Room-A can read the source-code of any Java bot in Room-B.

\section{ARCHITECTURAL PROCESS FLOW AND DATA LIFECYCLE}
Before describing how each of the individual lower level components will interact with one another, we need to define the overall Business Process Model and Notation (BPMN) as well as all possible sequential and/or parallel steps involved in the overall process of how data flows chronologically throughout the platform. Web based systems are generally built around traditional synchronous request-response cycles, whereas executing a 500-tick Tron match using compiled Java binaries will consume a great deal of computational resources.

Therefore, the BPMN design enforces an asynchronous workflow. After submitting an algorithm by a user, the front-end receives a non-descriptive "200 Accepted" response from the back-end within approximately 50ms. At this point, the payload has entered into a queuing buffer. A worker node then reads from the queuing buffer, compiles the binary code inside the execution sandbox, executes the game until completion, and records the JSON output to the database. Meanwhile, the client front-end continuously monitors its own background polling mechanism and waits passively for updates to be signaled regarding the status of the queuing buffer, and only fetches the heavier graphical rendering payload after receiving a clear signal from the server that the match is unambiguously terminated.
% PLACEHOLDER FOR BPMN DIAGRAM
\begin{figure}[hbp]
    \centering
    % A placeholder box that will render as a solid rectangle in the PDF
    % \rule{0.8\textwidth}{6cm}
    \includegraphics[width=0.9\textwidth]{pictures/BPMN.png}
    \caption{BPMN Diagram - Overall Platform Execution Workflow and Asynchronous Data Pipelining}
    \label{fig:bpmn_flow}
\end{figure}

\section{ACTOR USE CASE MODELLING}
The Use Case Diagram establishes a high-level boundary of the platform, explicitly mapping the authorized actions of various internal and external entities. The EMERGENT platform is driven by three distinct primary actors, each with explicitly segregated operational capabilities.

\subsection{The Administrator (Organizer)}
The Administrator holds the highest tier of localized platform privilege. Their primary goal is the successful facilitation of the tournament. The use cases constrained to this actor include:
\begin{itemize}
    \item \textbf{Room Provisioning:} Generating a new secure tournament room and obtaining the cryptographic PIN.
    \item \textbf{Tournament Configuration:} Defining the primary abstraction game (e.g., Tron or Connect4) that will be enforced for all players in that lobby.
    \item \textbf{Lobby Control:} Broadcasting the room code to participants and manually closing the room post-tournament to prevent further algorithmic submissions.
\end{itemize}

\subsection{The Participant (Student/Developer)}
The Participant interacts almost exclusively with the developer interface. Their actions are highly restricted strictly to the tenancy of the room they have explicitly joined.
\begin{itemize}
    \item \textbf{IDE Integration:} Selecting a designated compiler language, writing code within the Monaco IDE, and executing local syntax checks.
    \item \textbf{Code Submission:} Pushing versioned bots to the active matchmaker queue.
    \item \textbf{Test Match Invocation:} Requesting isolated, unranked execution sequences against static pre-calibrated AI Benchmark bots to test logic before committing entirely to the global ranked ladder.
    \item \textbf{Post-Action Review:} Browsing historical matches and utilizing the Canvas Visualizer to replay the exact deterministic progression of the executed match.
\end{itemize}

\subsection{The System Sandbox Actor (The Matchmaker)}
Unlike the human-driven actors, the Matchmaker is a continuous, automated system daemon. 
\begin{itemize}
    \item \textbf{Algorithmic Pairing:} Polling the active player roster and sorting them via a "Max-Information-Gain" scheduler to find the most statistically relevant pair.
    \item \textbf{Container Escalation:} Sending low-level execution demands to the Docker Daemon.
    \item \textbf{Rating Finalization:} Parsing the binary exit codes of the containers to recalculate and commit the comprehensive Glicko-2 participant rating shifts.
\end{itemize}

% PLACEHOLDER FOR USE CASE DIAGRAM
% \begin{figure}[hbp]
%     \centering
%     \rule{0.8\textwidth}{6cm}
%     \caption{UML Use Case Diagram depicting Admin, Participant, and Autonomous System Actor boundaries}
%     \label{fig:use_case}
% \end{figure}

\section{SYSTEM SEQUENCE DESIGN AND TIMING DIAGRAMS}
Sequence diagrams are essential for modeling the dynamic, inter-process communication behavior of the system, depicting the exact chronological flow of messages between object controllers over time. The fundamental operational pipeline in the EMERGENT platform is the "Submission to Execution Lifecycle", which requires strict transactional integrity to prevent corrupted database states or memory hanging.

The exact systemic sequence unfolds through the following chronological operations:
\begin{enumerate}
    \item The Participant's browser application parses the raw string arrays from the embedded Monaco editor. It bundles the payload alongside language configuration flags and invokes an asynchronous \texttt{multipart/form-data} POST request traversing the frontend routing layer.
    \item The Backend REST Controller validates the JWT or minimalist authentication headers, asserts that the active \texttt{room\_code} is still alive and accepting submissions, and securely flushes the bot script into a segregated local file system cache 


    (e.g., \texttt{/bots\_submissions/<room\_id>/<user>/}).
    \item Refusing to block the active thread, the API issues a "Pending Request Job" mapped to the Redis/Database queuing system and immediately responds to the HTTP connection with a \texttt{202 Accepted} response, restoring the client UI to a responsive, non-frozen state.
    \item Independently, the Background Matchmaker daemon thread executes its scheduler tick (typically every 1-2 seconds), sweeps the queue, and identifies the updated bot status. It queries the Glicko-2 tables to form an academically valid pairing.
    \item The internal Python Referee class is instantiated. It constructs an isolated, read-only Docker Volume referencing the newly saved files and dynamically provisions two ephemeral Docker containers (often utilizing the paused "Warm Pool" to eliminate cold-start lag). Both containers are connected strictly to each other and the host referee via standard IPC pipes.
    \item The primary turn-by-turn simulation loop commences. The Referee broadcasts the JSON formatted board state to \texttt{stdin} and waits up to a strict 1.0-second threshold to read the respective coordinates from \texttt{stdout}. The use of stdin and stdout makes it very easy for the participants to write their bots, focusing primarily on the logic and not on boilerplate code for handling input and output.
    \item Post-simulation (due to victory, algorithmic draw, timeout, or OOM exception), the Referee broadcasts a \texttt{SIGKILL} wiping the containers. The worker thread runs the complex Glicko-2 volatility iterations, maps the rating updates, and stringifies the dense \texttt{GameState.json} spatial payload.
    \item A final atomic commit pushes the aggregated results into the persistence layer. The Participant's frontend dashboard, acting via localized background polling, detects the updated schema timestamp and automatically invokes a React hook, injecting the new simulation log directly into the replay canvas.
\end{enumerate}

% PLACEHOLDER FOR SEQUENCE DIAGRAM
\begin{figure}[h]
    \centering
    % \rule{0.8\textwidth}{8cm}
    \includegraphics[width=0.9\textwidth]{pictures/sequence_project_1.pdf}
    \caption{UML Sequence Diagram delineating Code Submission mapping to Deterministic Container Execution}
    \label{fig:sequence_diagram}
\end{figure}

\section{SECURITY ARCHITECTURE AND SANDBOX ISOLATION}
A paramount design consideration for the platform involves threat mitigation. Designing an environment meant to execute completely untrusted remote code implies that malicious developers will attempt system escaping. 

The security architecture of EMERGENT is designed around "Zero-Trust Containerization." Instead of running code directly on the host's OS shell, the system utilizes Linux native namespaces and control groups (cgroups), orchestrated by the Docker Daemon, to create severely restricted micro-environments. 
\begin{itemize}
    \item \textbf{Network Isolation:} Containers are booted exclusively with the \texttt{--network=none} parameter. A bot theoretically infected with a reverse-shell payload cannot dial out to a remote command-and-control server. Furthermore, the two competing bots cannot communicate directly with one another over local sockets to collude; all data must strictly flow through the host Umpire structure via standard input/output pipes.
    \item \textbf{Process and Resource Throttling:} To heavily mitigate Denial-of-Service (DoS) and Fork-Bomb attacks (e.g., recursive child spawning), the execution environment design enforces rigid physical constraints. The environment utilizes cgroup limitations such as \texttt{--cpus="1.0"} (restricting execution to a single physical core) and \texttt{--memory="512m"}. In the event a script attempts to allocate arrays continuously, the Linux kernel acts as the ultimate arbiter, aggressively striking down the container process the exact millisecond memory usage breaches the threshold.
\end{itemize}

\section{RELATIONAL DATABASE DESIGN AND NORMALIZATION}
The structural pivot towards the "Room-First" MVP necessitated the creation of a highly normalized, multi-tenant relational schema. The primary goal of the SQL data layer is localized isolation; variables, leaderboard queries, and active algorithm paths must never encounter context spillage across separate, actively running tournaments.

The database operates primarily within Third Normal Form (3NF), ensuring data integrity while aggressively minimizing localized repetition. The decision to forgo massive ORMs in favor of direct query execution ensures that latency remains incredibly low even under the stress of rapid matchmaking commits.

\subsection{Core Entity-Relationship Mappings}
The Entity-Relationship (ER) model maps dependencies across four tightly interdependent canonical structures:

\begin{itemize}
    \item \textbf{Rooms Table (\texttt{rooms}):} This table functions as the primary namespace and isolation boundary for each independent game instance. It enforces globally unique \texttt{room\_code} values, typically generated via a high-entropy random or hash-based mechanism to prevent collisions and unauthorized access. In addition to identification, it stores critical configuration metadata, including the associated \texttt{game\_key}, which acts as a strict foreign constraint ensuring that only logically compatible agents (e.g., Tron bots vs Connect4 solvers) are evaluated within the same execution context. The table may also persist cryptographic artifacts such as hashed administrator credentials or control tokens, enabling secure room-level privilege management. This design guarantees both logical segregation and security across concurrent tournament environments.

    % PLACEHOLDER FOR ER DIAGRAM
\begin{figure}[h]
    \centering
    % \rule{0.8\textwidth}{7cm}
    \includegraphics[width=0.6\textwidth]{pictures/er_diagram_project_1.png}
    \caption{Structured Entity-Relationship (ER) Model displaying Third Normal Form and Room-First constraints}
    \label{fig:er_diagram}
\end{figure}

    \item \textbf{Participants Table (\texttt{participants}):} This entity represents the abstraction of a user within the scope of a specific room. It enforces a composite uniqueness constraint on (\texttt{room\_id}, \texttt{display\_name}), ensuring that identity is strictly scoped locally—thereby preventing namespace collisions across rooms. For instance, a participant named ``Alex'' in Room A is treated as a completely independent entity from ``Alex'' in Room B, 
    
    with no shared relational linkage. Beyond identity, this table stores references to the participant's submitted bot logic, typically as filesystem paths or object storage URIs pointing to serialized binaries or compiled artifacts. Additional metadata such as submission timestamps, version hashes, or execution language identifiers may also be included to support reproducibility, auditing, and sandboxed execution workflows.

    \item \textbf{Ratings Metadata Table (\texttt{room\_ratings}):} This table is intentionally decoupled from the core participant schema to optimize for high-frequency, write-intensive operations driven by continuous match execution. It encapsulates the full state required by the rating algorithm (e.g., a Glicko-2 or similar system), including the scalar \texttt{rating} representing skill estimation, the \texttt{rd} (Rating Deviation) quantifying uncertainty or confidence intervals, and the \texttt{vol} (Volatility) capturing the degree of performance fluctuation over time. By isolating this data, the system minimizes lock contention and enables efficient batch or incremental updates without impacting participant identity records. This separation also facilitates potential future extensions, such as maintaining historical rating snapshots or supporting multiple rating systems in parallel.

    \item \textbf{Transaction Table (\texttt{room\_matches}):} The above table is a source for storing transactional information regarding every competitive interaction that occurs in a room. For this reason, it creates a many to many relationship between two different participants by way of \texttt{participant\_a\_id} and \texttt{participant\_b\_id}. In addition, the table captures each individual match. Furthermore, the schema captures the outcome of each match by 
    
    using strongly-typed categorical enum types. These allow for an extremely high degree of granularity when determining the outcome of a match. As well, they can be used to determine whether a participant won due to their strategy or if the participant experienced a run-time exception (for example segmentation fault or syntax error) or was forced to terminate due to one of the standard I/O time-outs. Finally, the table contains large amounts of serialized data-blobs, often stored in JSON format. Each blob represents the complete state of the entire match replay (that is, 2-dimensional coordinates, number of game ticks, or event logs). Therefore, the table acts as a means for performing deterministic replays of matches as well as a mechanism for debugging and analyzing how agents behaved after the fact. As such, the table serves as both a transactional record and a store for artifacts related to diagnosis.
\end{itemize}



% ====================
% CHAPTER 4: SYSTEM IMPLEMENTATION
% ====================
\chapter{SYSTEM IMPLEMENTATION}

System Implementation is the critical phase where the theoretical system designs, use cases, and database schemas are translated into functional source code. For EMERGENT, this phase required integrating diverse technological stacks: a modern React front-end, a robust Flask python backend, and a low-level OS-dependent Docker execution engine. This chapter provides a highly detailed breakdown of how each core module was engineered, focusing heavily on the algorithmic logic, pipeline management, and technical solutions to edge-case anomalies.

\section{FRONTEND UI AND WORKSPACE IMPLEMENTATION}
The user interface serves as the primary conduit for participant interaction. Built using React.js and scaffolded with Vite to ensure blazing-fast Hot-Module Replacement (HMR) during the development lifecycle, the frontend is engineered as a Single Page Application (SPA). This SPA architecture prevents jarring page reloads, which is critical when users are actively coding and monitoring live match queues.

\subsection{Monaco Editor Workspace Integration}
At the heart of the developer experience is the integrated web IDE workspace. Rather than relying on basic HTML text areas, the platform embeds the Monaco Editor—the exact declarative text-editing component that powers Microsoft's Visual Studio Code.
\begin{itemize}
    \item \textbf{Language Agnostic Syntax Highlighting:} The Monaco instance is dynamically configured to support multiple compiler languages. When a user selects Python, Java, or C/C++ from the dropdown, the editor swaps its internal lexer rules, providing immediate keyword highlighting, scope folding, and intelligent bracket matching.
    \item \textbf{State Management:} As the user types, the code string is maintained in the localized React Context state. Upon clicking the "Submit" module, the frontend packages this raw string, along with the user's explicit \texttt{room\_code} and \texttt{display\_name}, into a multi-part form data structure to be transmitted via the REST HTTP protocol to the backend.
\end{itemize}

\subsection{Modular Replay Visualizer}
Because EMERGENT supports multiple abstraction games (e.g., Tron, Connect4), the Visualizer component is designed as a dynamic schema consumer rather than hardcoding rendering logic. 
\begin{itemize}
    \item \textbf{Data Ingestion:} Once a match concludes, the client fetches the \texttt{GameState.json} payload from the Flask API. This payload outlines every grid coordinate, bounding box, and frame variation generated by the Referee Engine.
    \item \textbf{Frame-by-Frame Rendering:} Using standard HTML5 Canvas and advanced CSS Grid layouts, the Visualizer iterates through the array of game states. It exposes a "Flight Recorder" control panel to the user, complete with Play, Pause, Next Tick, and Previous Tick buttons. By decoupling the render logic from the game logic, new topological games can be displayed with minimal frontend refactoring.
\end{itemize}

\section{BACKEND ORCHESTRATION AND APIs}
The main server application acts as the synchronization hub. Developed in Python utilizing the Flask micro-framework, it maps URL routes to underlying Python controller functions.

\subsection{Room Management and Minimalist Authentication}
To prioritize rapid tournament spooling for LAN clubs and academia, heavy OAuth integrations were discarded in favor of a "Minimalist Authentication" model. 
When an Administrator hits the \texttt{POST /api/rooms} endpoint, the server executes a randomized cryptographic generator to spit out a unique 6-character alphanumeric room code and a 4-digit admin PIN.
\begin{itemize}
    \item \textbf{Participant Binding:} When a student submits their code, the server queries the database. If the \texttt{(room\_code, display\_name)} composite payload does not exist, it registers a new participant and hashes that initial submit password utilizing SHA-256 with an aggressive salt. 
    \item \textbf{Subsequent Authorization:} If a student attempts to update their bot or request a benchmark test, the middleware compares the provided password payload against the stored hash. This completely sandboxes user profiles to specific rooms.
\end{itemize}

\section{MATCHMAKING AND RATING ENGINE}
To prevent server I/O blocking during heavy compilations, the system abandons manual matchmaking in favour of an autonomous Background Worker and live scheduler architecture. 

\subsection{Live Room Scheduler and Max Information Gain}
The fundamental problem in continuous automated tournaments is deciding \textit{who plays whom} and \textit{when}. Random pairings are highly inefficient; pitting a beginner against a multi-tournament Grandmaster yields almost no new statistical data about either player's true skill. To resolve this, EMERGENT implements a Live Room Scheduler Tick policy that utilizes a "Max-Information-Gain" pairing algorithm.

The scheduler runs periodic ticks, evaluating the entire pool of participants within a room based on three discrete phases:

\begin{enumerate}
    \item \textbf{Calibration Phase (High RD):} When a user submits a bot for the first time, their Rating Deviation (RD) defaults to 350. The scheduler prioritizes pairing this new bot against standardized "Tiered Benchmark Bots" (pre-calibrated internal AI) or veteran players with very low RDs. This rapidly brackets the new user's skill, immediately pulling their RD down and preventing leaderboards from becoming polluted with uncalibrated anomalies.
    \item \textbf{General Tick Policy \& Max Information Gain:} Once calibrated, the scheduler computes the relative distance between all active bots. It pairs bots where the predicted outcome is closest to 50/50 (i.e., their ratings are nearly identical). A match between evenly matched opponents yields the "Maximum Information Gain" mathematically, causing the greatest meaningful adjustment to their ratings post-match.
    \item \textbf{Convergence Stop Gate \& Rematch Caps:} To prevent server exhaustion and stop two bots from uselessly farming each other, the system tracks bot-version rematches. A hard configurable rematch cap restricts identical bot versions from fighting repeatedly. Furthermore, a Room Convergence Stop Gate is implemented; if all bots in a room have played each other and their RD falls below a stable threshold (e.g., RD $<$ 50), the scheduler safely pauses matchmaking for that room until a user submits new, updated code.
\end{enumerate}

\subsection{Glicko-2 Algorithmic Implementation}
Upon the successful completion of a scheduled match, the worker thread intercepts the Win (1), Loss (0), or Draw (0.5) termination code. Standard Elo algorithms are fundamentally flawed for bot tournaments because they naively assume a player's current rating is absolutely precise. EMERGENT instead implements Mark Glickman's \textbf{Glicko-2 System}, which introduces a multidimensional measure of skill. 

Each participant's profile tracks three dynamic variables:
\begin{itemize}
    \item \textbf{Rating ($r$):} The raw measure of absolute skill. Default starting value is 1500.
    \item \textbf{Rating Deviation ($RD$):} Represents the mathematical uncertainty of the rating. Default starting value is 350. As a player completes more matches, the system becomes more confident in their skill, and RD approaches 0.
    \item \textbf{Rating Volatility ($\sigma$):} Measures the degree of expected statistical fluctuation in a player's rating. A highly consistent player will have a low volatility, whereas a highly erratic bot will have a high volatility. Default starting value is 0.06.
\end{itemize}

\subsubsection*{Mathematical Computation Logic}
When processing the result, the engine first converts the base Glicko ratings into the standard Glicko-2 scale parameters using the scale multiplier $173.7178$:
\begin{equation} \label{eq:scale}
    \mu = \frac{r - 1500}{173.7178} \quad \text{and} \quad \phi = \frac{RD}{173.7178}
\end{equation}
Where in Equation \ref{eq:scale}:
\begin{itemize}
    \item $\mu$ is the converted Glicko-2 rating parameter.
    \item $\phi$ is the converted Glicko-2 rating deviation parameter.
    \item $r$ is the participant's original base Rating.
    \item $RD$ is the participant's original base Rating Deviation.
\end{itemize}

To compute the impact of an opponent's rating deviation on the match outcome, the algorithm applies the weighting function:
\begin{equation} \label{eq:weight}
    g(\phi) = \frac{1}{\sqrt{1 + \frac{3\phi^2}{\pi^2}}}
\end{equation}
Where in Equation \ref{eq:weight}:
\begin{itemize}
    \item $g(\phi)$ represents the weighting function that diminishes the impact of a rating change if the opponent's rating is highly uncertain.
    \item $\phi$ is the opponent's rating deviation on the Glicko-2 scale.
    \item $\pi$ is the mathematical constant Pi.
\end{itemize}

This allows the engine to calculate the expected outcome of a match between the evaluated player and the opponent:
\begin{equation} \label{eq:expected}
    E(\mu, \mu_j, \phi_j) = \frac{1}{1 + \exp(-g(\phi_j)(\mu - \mu_j))}
\end{equation}
Where in Equation \ref{eq:expected}:
\begin{itemize}
    \item $E$ is the expected probability of the evaluated player winning the match.
    \item $\mu$ is the evaluated player's converted Glicko-2 rating.
    \item $\mu_j$ is the specific opponent's converted Glicko-2 rating.
    \item $\phi_j$ is the specific opponent's converted rating deviation.
\end{itemize}

Once expectations are established, the engine evaluates the estimated variance of the player's rating based on the collected match outcomes alongside the estimated scalar improvement:
\begin{equation} \label{eq:variance}
    v = \left[ \sum_{j=1}^{m} g(\phi_j)^2 E(1 - E) \right]^{-1}
\end{equation}
\begin{equation} \label{eq:improvement}
    \Delta = v \sum_{j=1}^{m} g(\phi_j)(s_j - E)
\end{equation}
Where in Equations \ref{eq:variance} and \ref{eq:improvement}:
\begin{itemize}
    \item $v$ is the estimated variance of the player's rating resulting from the rating period.
    \item $\Delta$ is the estimated improvement in the true rating based on the match results.
    \item $m$ is the total number of opponents played.
    \item $j$ denotes the index of a specific opponent in the summation.
    \item $s_j$ represents the objective, observable match score against opponent $j$ ($1$ for win, $0.5$ for draw, $0$ for loss).
\end{itemize}

Finally, the backend utilizes an iterative numerical procedure (specifically utilizing the Illinois algorithm) to find the new Volatility ($\sigma'$), which is plugged back into the primary equations to compute the post-match $\phi'$ and $\mu'$. These values are instantly reverted back to the standard $r$ and $RD$ scales and atomically committed back into the SQLite \texttt{room\_ratings} table. 

This rigorous mathematical integration drastically enhances algorithmic integrity, heavily rewarding users who consistently write robust strategies while aggressively demoting architectures that rely on lucky edge-case victories.

\section{THE DETERMINISTIC DOCKER SANDBOX}
The most complex engineering challenge within the EMERGENT stack is the Referee Engine. It must execute highly untrusted, possibly malicious, and inevitably buggy code provided by students without compromising the host architecture.

\subsection{Sandboxing and Ephemerality}
Every single match invokes the Docker daemon. The referee spins up dedicated, highly restricted containers devoid of internet routing access via the \texttt{--network=none} flag. Furthermore, strict quota limits are heavily enforced directly onto the Docker API using flags such as \texttt{--memory="512m"} and \texttt{--cpus="1.0"}. If a student writes a massive memory leak (e.g., an infinite appending array), the container process is forcefully killed by the OS kernal (yielding an Out-Of-Memory OutOfMemory OOM exception), which the Referee catches, interprets, and communicates cleanly back to the client as an "Internal Crash Loss."

\subsection{Interprocess Communication (IPC) and Anomaly Handling}
The Referee Engine utilizes Python's built-in \texttt{subprocess.Popen} module to pipe streams directly into the running Docker containers. At the start of a turn, the engine writes the physical board state to standard input (stdin) of the student's bot, and heavily monitors standard output (stdout) for the return coordinates.

Significant technical implementations were required to handle I/O anomalies:
\begin{itemize}
    \item \textbf{Non-terminating Strings:} If a user's code attempts to output a move but fails to flush a newline character (\texttt{\textbackslash n}), typical buffered pipeline reading will halt indefinitely. The system bypasses this by forcing select-based asynchronous reads with strict millisecond timeouts.
    \item \textbf{Standard Error Blocking:} Normally, if a bot throws a syntax exception and writes heavily to \texttt{stderr}, it can block the buffer and crash the host thread. EMERGENT spins off secondary daemon threads to explicitly drain \texttt{stderr} pipes concurrently, ensuring the worker node gracefully identifies script crashes and immediately awards the opposing bot the win.
\end{itemize}

\subsection{Game Plugin Registry}
To ensure the execution engine is decoupled from any hardcoded game rules, all game logic is routed through a \texttt{bases.py} interface. A "Tron" match is executed by loading the \texttt{TronPlugin}, which acts purely as an umpire—it verifies if an emitted coordinate is valid, computes collisions, and appends the tick state to the evolving \texttt{GameState.json} ledger.

\section{IMPLEMENTATION SCREENSHOTS}

The translation of the architectural design into a fully functional product is best demonstrated through the graphical user interface. The following subsections detail the finalized visual components of the EMERGENT platform, showcasing how the underlying React.js Single Page Application manages complex data workflows while providing a frictionless experience for the participants.

\subsection{Platform Entry and Authentication}
The initial interaction point for developers is the landing view. Reflecting the "Room-First" architecture, this portal serves as the primary access gate. Participants are prompted to enter an active Room Code securely generated by the Umpire, alongside their chosen display name. The minimalist design completely foregoes complex OAuth flows or registration portals, prioritizing rapid onboarding suitable for localized hackathons and academic environments.

% PLACEHOLDER FOR SCREENSHOTS
\begin{figure}[hbp]
    \centering
    % \rule{0.8\textwidth}{5cm}
    \includegraphics[width=0.9\textwidth]{pictures/2026-03-30_00-11.png}
    \caption{Implementation Screenshot - Single Page Application Landing View}
    \label{fig:screen_landing}
\end{figure}

\subsection{Integrated Developer Workspace}
Once authenticated, the participant is transitioned into the core development workspace. Figure \ref{fig:screen_workspace} illustrates the seamless integration of the Monaco Editor within the browser. The interface immediately provides intelligent syntax highlighting, code folding, and line numbering—eliminating the need for any local IDE configuration. Adjacent to the editor, real-time compilation and execution logs stream directly from the backend, providing immediate textual feedback upon code submission.

\begin{figure}[H]
    \centering
    % \rule{0.8\textwidth}{5cm}
    \includegraphics[width=0.9\textwidth]{pictures/2026-03-30_00-05.png}
    \caption{Implementation Screenshot - Web Workspace demonstrating Monaco IDE Integration and Execution Logs}
    \label{fig:screen_workspace}
\end{figure}

\subsection{Deterministic Replay Visualization}
To resolve the "Black Box" debugging problem, the platform reconstructs the match history visually. Figure \ref{fig:screen_replay} displays the Frame-by-Frame Visual Replay System. Leveraging HTML5 Canvas logic to parse the dense \texttt{GameState.json} artifacts, the UI paints a clear, colorful 2D abstraction of the game arena (in this case, the Tron lightcycle grid). The intuitive playback controls located on the interface allow the developer to pause, rewind, and scrutinize exact turn maneuvers, isolating strategic flaws with chronological precision.

\begin{figure}[h]
    \centering
    % \rule{0.8\textwidth}{5cm}
    \includegraphics[width=0.9\textwidth]{pictures/2026-03-30_00-04.png}
    \caption{Implementation Screenshot - Frame-by-Frame Visual Replay System (Tron Representation)}
    \label{fig:screen_replay}
\end{figure}

\subsection{Real-Time Statistical Leaderboard}
Finally, the dynamic assessment goals of the platform are realized through the live leaderboard dashboard, depicted in Figure \ref{fig:screen_leaderboard}. As the background worker pool constantly resolves queued matches, this panel dynamically populates without requiring a page refresh. It exposes the complex mathematical metrics computed by the Umpire, clearly displaying each participant's current absolute Glicko-2 Rating ($r$), their Rating Deviation ($RD$), and cumulative Win/Loss records. Ensuring that participants can constantly evaluate their strategic performance relative to their peers.

\begin{figure}[h]
    \centering
    % \rule{0.8\textwidth}{5cm}
    \includegraphics[width=0.9\textwidth]{pictures/2026-03-30_00-34.png}
    \caption{Implementation Screenshot - Real-Time Dashboard and Glicko-2 Leaderboard view}
    \label{fig:screen_leaderboard}
\end{figure}

% ====================
% CHAPTER 5: TESTING
% ====================
\chapter{TESTING}

A central component in the Systems Development Life Cycle (SDLC) is software testing. The purpose of software testing is to verify that a software platform has been developed correctly and will provide stable and secure functionality for all users prior to deployment. Due to its nature, EMERGENT requires more than just normal quality assurance testing through dynamic testing; rather, due to the untrusted and possibly malicious user-provided code executed, it requires an additional layer of dynamic testing as a critical security requirement.

The comprehensive testing strategies employed for the development of this system include unit testing of individual components or modules of the system (i.e., mathematical formulas), as well as extensive edge-case testing at multiple levels to evaluate the robustness of the Docker sandbox against various types of system level attacks such as resource depletion and pipe blocking.

\section{TESTING STRATEGY AND METHODOLOGY}
The testing strategy employed for the EMERGENT architecture utilizes a hybrid approach, combining both White-Box testing (where the internal structure and implementation are known to the tester) and Black-Box testing (where scripts strictly evaluate the input-output behaviors without architectural awareness).

Given the modular nature of the application, the testing lifecycle was broken down into the following hierarchical tiers:
\begin{enumerate}
    \item \textbf{Unit Testing:} Validating individual atomic functions (such as the Glicko-2 module and Password Hashing algorithms) in total isolation.
    \item \textbf{Integration Testing:} Ensuring the localized components (e.g., the Flask routing controllers interfacing with the SQLite database) successfully communicate.
    \item \textbf{Sandbox Resilience Testing:} A highly specialized testing tier unique to EMERGENT, wherein malicious or malformed bot scripts are purposefully injected to attempt to break the server host.
    \item \textbf{System and End-to-End (E2E) Testing:} Simulating a complete user journey—from interacting with the React/Monaco frontend, submitting code, queueing the match, and rendering the final JSON replay.
\end{enumerate}

\section{UNIT TESTING}
Unit tests were constructed to verify the operational logic of the backend application utilizing Python’s built-in \texttt{unittest} and \texttt{pytest} frameworks. By heavily mocking the database connections, these tests execute in milliseconds, allowing rapid iterative verification.

\subsection{Glicko-2 Algorithmic Verification}
The most critical unit tests involve the matchmaking rating system. Because algorithmic integrity is absolute, scripts were written to inject predetermined arrays of matches to assert the mathematical outputs.
\begin{itemize}
    \item \textbf{Expected Volatility Shift:} Scripts verified that a standard 1500 Rating ($r$) player who suffered a catastrophic loss to an 800 Rating player experienced an intense spike in their Rating Volatility ($\sigma$) parameter, proving the Illinois algorithm's iteration bounds.
    \item \textbf{RD Reduction:} The system asserted that regardless of win or loss, the execution of 5 consecutive matches successfully reduced a new player's Rating Deviation (RD) from $350$ down to the expected baseline of circa $120$.
\end{itemize}

\subsection{Cryptographic Hash Validation}
Within \texttt{test\_backend.py}, unit handlers were established corresponding to the minimalist authentication logic mapping. The tests verify that the 
\texttt{generate\_password\_hash} algorithm successfully binds a randomized salt to every newly created participant, and that identical raw string inputs correctly resolve to successful boolean true during the 
\texttt{check\_password\_hash} routine.

\section{INTEGRATION TESTING}
Integration testing mapped the transactional boundaries between the routing layer, the database layer, and the asynchronous matchmaker. The testing script \texttt{test\_room\_flow.py} (located within the \texttt{scripts/active/} directory) automated the invocation of these heavily interdependent modules.

\subsection{Room Schema Isolation}
A fundamental paradigm of the MVP architecture is the "Room-First" philosophy. To guarantee that independent tournaments do not undergo cross-contamination, integration tests were run to simulate concurrent lobbying.
\begin{itemize}
    \item \textbf{Execution Flow:} The script fires \texttt{POST /api/rooms} twice to generate "Room A" and "Room B". User 'Alice' is appended to Room A, and an autonomous match is executed changing her rating to 1550.
    \item \textbf{Assertion Check:} A secondary test script attempts to query Alice's profile utilizing Room B's contextual headers. The API must firmly reject this via a HTTP 404 (Not Found) or 401 (Unauthorized), affirmatively proving strict database tenancy isolation.
\end{itemize}

\section{SANDBOX RESILIENCE AND LOGIC TESTING (EDGE CASES)}
Standard web applications rarely have to account for malicious user inputs executing arbitrary binaries natively. EMERGENT’s execution engine was subjected to aggressive "Chaos Testing" to expose critical vulnerabilities regarding container escape, OS thread locking, and algorithmic boundary breaking. In accordance with the design specifications, the following tests were conducted specifically analyzing the \textbf{Performed Test}, the \textbf{Observed Result}, and the \textbf{Correction Applied}.

\subsection{Resource Exhaustion and OOM Exceptions}
\textbf{The Performed Test:} A participant submits a bot coded to infinitely append integers to a global array structure in memory (e.g., \texttt{while True: arr.append(1)}). The test evaluates whether a single bot can consume all available server RAM and orchestrate a Denial-of-Service (DoS) attack. \\
\textbf{The Observed Result:} During initial alpha testing, the host machine locked up entirely as the container consumed 100\% of system resources, requiring a hard server reboot. \\
\textbf{The Correction Applied:} The Docker provisioning module was updated to enforce strict \texttt{cgroup} limitations, explicitly passing the \texttt{--memory="512m"} flag to the container. Upon re-testing, the Linux kernel forcefully executed a SIGKILL (Signal 9) on the container process the exact millisecond it breached 512MB. Crucially, the string parsing logic was corrected so the Referee engine elegantly catches the \texttt{Out-Of-Memory (OOM)} exception code instead of crashing, dynamically generating a JSON payload indicating an "Internal Crash Loss" for that participant.

\subsection{Standard Error (stderr) Pipe Buffering Anomaly}
\textbf{The Performed Test:} A "noisy" bot deliberately structured to vomit 10 Megabytes of garbage crash logs to \texttt{sys.stderr} was tested to evaluate OS-level interprocess communication limits. \\
\textbf{The Observed Result:} The Python \texttt{subprocess.Popen} module connects to a spawned container via pipes. The default OS pipe buffer capacity is merely 64KB. When the buffer filled completely, the OS blocked the container from writing further. Concurrently, the Host Referee, blindly waiting to read from \texttt{stdout}, blocked indefinitely. This resulted in a permanent "Thread Freeze"—creating zombie worker threads that completely halted the matchmaking queue. \\
\textbf{The Correction Applied:} The core Referee architecture was entirely rewritten to spawn independent, asynchronous daemon threads dedicated exclusively to draining the \texttt{stderr} pipes concurrently. Following this correction, the test proved that the worker node could successfully identify the bot crash instantly without locking the system buffer.

\subsection{Non-Terminating String IO (The Newline Bug)}
\textbf{The Performed Test:} The system evaluated how the worker node handles incorrectly formatted I/O string flushes. The test bot script executed: \texttt{print("UP", end="")} without a native trailing newline. \\
\textbf{The Observed Result:} The baseline Referee's \texttt{readline()} loop blocked indefinitely waiting for an EOF or newline delimiter. The match hung until the global container timeout stepped in, meaning the system inaccurately logged a severe logic string mismatch simply as a generic "Timeout Delay". \\
\textbf{The Correction Applied:} Custom \texttt{select()} OS-level hardware pollings were integrated into the \texttt{readline} engine. The correction ensures the buffer physically checks for the presence of standard ASCII streams independently, and if malformed output is caught, it correctly labels the failure as an "I/O String Mismatch Exception" rather than a hanging timeout.

\subsection{Grid Dimensionality and Boundary Frame Breaking}
\textbf{The Performed Test:} Early testing evaluated what would occur if a bot assumed the Tron arena was a hardcoded $15 \times 15$ grid, and the platform organizer dynamically altered the match parameters through the admin panel to a larger $30 \times 30$ grid boundary. \\
\textbf{The Observed Result:} The bots physically stopped playing matches. As they traversed past index 15, the frontend React Visualizer scaling formula broke, attempting to draw coordinate lines outside the HTML5 canvas boundaries. Simultaneously, the bots crashed throwing \texttt{IndexOutOfBounds} array exceptions because they were unaware of the active window size. \\
\textbf{The Correction Applied:} A major systemic correction was implemented to standardize the \texttt{GameState} payload schema. The backend Umpire logic was corrected to explicitly broadcast a \texttt{grid\_size} JSON header to every bot's \texttt{stdin} on standard initialization tick (Tick 0). Additionally, the frontend Canvas rendering equations were rewritten to dynamically calculate their scalar view-port (Window Size) by hooking immediately into this header, allowing flawless play regardless of dynamic grid sizes.

\subsection{Matchmaking Calibration (The High RD Tick Anomaly)}
\textbf{The Performed Test:} Evaluated the trajectory of a brand new bot being injected heavily into a fully saturated queue of 50 existing bots to monitor how the Elo rating scales. \\
\textbf{The Observed Result:} Without calibration-specific routing, the Live Matchmaker repeatedly threw the new bot (with a default 1500 Rating and 350 RD) into random queues. This led to highly unbalanced matches where the uncalibrated bot destroyed the integrity of the established leaderboard because it yielded extremely poor "Information Gain" for veteran bots playing it. \\
\textbf{The Correction Applied:} A "Calibration Isolation Tick" logic gate was structurally coded into the Redis queue scheduler. The correction intercepts any bot possessing an $RD > 200$ and forcefully prevents it from playing the general player population. Instead, the Matchmaker routes it exclusively against static Pre-Calculated Benchmark tier bots. This rapidly drops the RD variance parameter within five matches, yielding mathematically stable general queue conditions.

\section{END-TO-END (E2E) AND SYSTEM TESTING}
E2E testing replicates the absolute client experience. Unlike isolated scripts, these tests run fully across the React.js GUI.

\subsection{Monaco Editor and Frontend Evaluation}
Tests were conducted to evaluate that a user could seamlessly switch contexts. It verified that changing the dropdown from \texttt{Python} to \texttt{Java} accurately updated the Monaco Editor’s Language Services provider, and that issuing a "Submit" event correctly packaged the multipart \texttt{bot\_zip\_file} payload through the modern Fetch API boundaries cleanly.

\subsection{Replay Rendering Accuracy}
The Visualizer was tested using a heavily populated \texttt{GameState.json} dummy file containing over 500 game ticks. The system passed the visual assertion test by smoothly running the Canvas rendering at 60 Frames Per Second (FPS) without resulting in noticeable DOM lag, effectively proving the decoupling of the Game Logic implementation from the React rendering components.

\section{FORMAL TEST CASE SPECIFICATIONS}
The following tables formally document a cross-section of the explicit system test cases executed to guarantee compliance with the Software Requirement Specifications.

\vspace{0.5cm}
\begin{table}[ht]
\centering
\caption{System Implementation Test Cases - Module 1: API and Orchestration}
\label{tab:test_cases_api}
\vspace{0.2cm}
\renewcommand{\arraystretch}{2.0} % Stretches the rows making the table very large and readable
% \begin{tabular}{|p{1.5cm}|p{5cm}|p{4cm}|p{4cm}|p{1.5cm}|}
\begin{tabular}{|p{1.5cm}|p{4.2cm}|p{3.8cm}|p{3.8cm}|p{1.4cm}|}
\hline
\textbf{Test ID} & \textbf{Test Description} & \textbf{Expected Output} & \textbf{Actual Output} & \textbf{Status} \\ \hline
\texttt{TC-01} & Submit Python code with valid syntax and correct \texttt{room\_code}. & Code accepted, queued, HTTP 200 OK. & Code accepted, pending job created in DB, 200 OK returned. & \textbf{PASS} \\ \hline
\texttt{TC-02} & Submit Java code missing the compilation script (\texttt{run.sh}). & Referee catches missing artefact, returns Error. & \texttt{Validation Exception} raised, frontend displays error envelope. & \textbf{PASS} \\ \hline
\texttt{TC-03} & Create a room using the Admin Panel. & System generates exact 6-char code and 4-digit PIN. & 6-char UUID and hashed PIN successfully written to SQLite. & \textbf{PASS} \\ \hline
\texttt{TC-04} & Send unauthorized test match request with wrong password hash. & HTTP 401 Unauthorized middleware rejection. & 401 Unauthorized securely returned without touching queue. & \textbf{PASS} \\ \hline
\end{tabular}
\end{table}

\vspace{0.5cm}

To validate both the functionality and safety of the Docker-based execution sandbox in addition to confirming its reliability and accuracy (i.e., validating the API layer), a rigorous testing strategy must be implemented to test for both the security properties of this component as well as the isolation guarantees of this sandbox, as well as its fault-tolerance.

This sandbox will execute user-supplied code within an isolated environment; often under uncertain and possibly hostile/ malicious input conditions, and at times with very high concurrency levels (e.g. live tournament conditions). Therefore, the testing strategy for this component goes beyond verifying its function and includes both stress-testing and robustness-testing of the execution environment itself.

Below is a summary table of some examples of focused stress-tests designed to determine how the sandbox behaves under various load and boundary condition scenarios. These tests were developed to ensure that when executing code within the sandbox environment, all resource limits are properly enforced; the appropriate execution boundaries remain intact; and, regardless if the submitted code was either mangled/badly formatted, inefficient or malicious in intent, none of these can/will cause the integrity of the host system or the orchestrator layer to be compromised. Specific emphasis is placed upon process isolation, restriction of file-system accesses from processes executing in the sandbox, proper enforcement of timeouts, and successful cleanup of containerized workload(s) after execution has been completed.

\begin{table}[H]
\centering
\caption{System Implementation Test Cases - Module 2: The Execution Sandbox}
\label{tab:test_cases_sandbox}
\vspace{0.2cm}
\renewcommand{\arraystretch}{2.0}
% \begin{tabular}{|p{1.5cm}|p{5cm}|p{4cm}|p{4cm}|p{1.5cm}|}
\begin{tabular}{|p{1.5cm}|p{4.2cm}|p{3.8cm}|p{3.8cm}|p{1.4cm}|}
\hline
\textbf{Test ID} & \textbf{Test Description} & \textbf{Expected Output} & \textbf{Actual Output} & \textbf{Status} \\ \hline
\texttt{TC-05} & Inject an infinite \texttt{while} loop into bot script. & Referee enforces absolute 1.0s turn timer. & Container halted after exactly 1.03s. Timeout Loss recorded. & \textbf{PASS} \\ \hline
\texttt{TC-06} & Inject huge unallocated memory array into Python list. & OS kernel issues SIGKILL when exceeding 512MB RAM. & Docker catches memory spike. OOM exception properly parsed by Worker. & \textbf{PASS} \\ \hline
\texttt{TC-07} & Run Benchmark Tier 1 Bot against itself. & Bit-for-Bit Deterministic JSON output generated. & Identical seed generated identical replay byte-for-byte. & \textbf{PASS} \\ \hline
\texttt{TC-08} & Submit 50 concurrent matches to Redis Queue. & Host processes gradually via Warm Pool without crash. & Queue populated and gracefully exhausted across 3 nodes over 20s. & \textbf{PASS} \\ \hline
\end{tabular}
\end{table}
The system has shown it can execute a variety of execution scenarios through successful completion of the above mentioned test cases. The various scenarios include but are not limited to chaotic logic patterns and other types of adversarial input. Execution is demonstrated as being dynamic and isolated within the sandbox environment and therefore does not leak resources nor expose vulnerability to the Host Orchestrator. Therefore the host orchestrator will remain secure and reliable, under both high load and hostile user activity, ensuring the reliability and security of the Platform when deployed in production environments.

% By successfully passing the aforementioned test cases, the execution sandbox proves capable of handling chaotic logic and adversarial code submission dynamically, securing the host orchestrator from malicious interference.

% ====================
% CHAPTER 6: RESULTS AND DISCUSSION
% ====================
\chapter{RESULTS AND DISCUSSION}

\section{SYSTEM OUTPUTS}
The EMERGENT platform successfully generates several distinct outputs that formulate the complete user experience. By transitioning from a standard terminal interface to a robust web-based graphical user interface (GUI), the system outputs were specifically designed to provide immediate, actionable feedback to the developers.

\subsection{Interactive Web Workspace and Real-Time Feedback}
The primary output for the user is the Monaco-powered IDE, perfectly mirroring professional environments like Microsoft Visual Studio Code. It successfully outputs real-time syntax checking and language-specific compilation environments without requiring any local setups. By integrating declarative Abstract Syntax Tree (AST) parsing directly into the browser, the workspace provides an immediate cognitive loop where developers receive compile-time errors within milliseconds of pressing the submit button. This instant diagnostic output heavily reduces the friction traditionally associated with configuring localized development environments.

\subsection{Comprehensive Architectural Logs (\texttt{GameState.json})}
Instead of simply outputting a binary 'Win' or 'Loss' state upon the conclusion of a match, the execution engine generates a massive, tightly structured JSON payload. This payload serves as a comprehensive architectural log, encompassing every single move, grid coordinate selection, wall collision event, and resource allocation occurring per game tick. By capturing the complete lifecycle of the matched algorithms, the JSON output acts as a deterministic ledger that proves exactly why an outcome was reached, completely resolving the ambiguity present in traditional pass/fail testing platforms.

\subsection{The Replay Visualizer and Canvas Engine}
The React frontend ingests the aforementioned JSON artifact and outputs a two-dimensional HTML5 Canvas rendering. This "Flight Recorder" interface outputs a buttery-smooth 60 Frames Per Second (FPS) playback. It enables engineers to scrub through their algorithm's exact decisions frame-by-frame, pause at critical junctures, and rewind to analyze the precise tick where their computational logic deviated from their intended strategy. The separation of game logic from rendering logic ensures that the frontend maintains high execution speeds regardless of the complexity of the underlying algorithm.

\subsection{Dynamic Leaderboards and Participant Metrics}
The system continuously outputs live-updating leaderboards utilizing the Glicko-2 ratings parameterization ($\mu$, $\phi$, $\sigma$). This output grants server administrators and tournament organizers an immediate top-down view of algorithmic dominance across the LAN tournament. By displaying the Rating Deviation ($RD$) prominently alongside the overall score, the leaderboard introduces absolute transparency, allowing participants to visually understand whether their rank is derived from statistical certainty or mere initial calibration fluctuation.

\section{PERFORMANCE EVALUATION AND RESULTS}
Through aggressive deployment testing, the architecture yielded exceptional performance results, proving its viability for live LAN-party contexts where latency and resource distribution are highly scrutinized. 

\subsection{Latency, Throughput, and The Warm Pool Advantage}
By abandoning the synchronous web-server execution model, the API routing layer maintained an average response latency of strictly under 50 milliseconds for code submissions. The critical bottleneck in containerized platforms is often the startup latency incurred when spinning up virtualized environments. By implementing a Docker "Warm Pool"—a technique where baseline environments are pre-initialized and held in a paused state—the platform drastically reduced the container cold-start delay. Empirical test results showed match initiation times dropping from an average of over 2.5 seconds down to roughly 400 milliseconds. This optimization allowed the background asynchronous queue to chew through heavy match backlogs seamlessly, processing up to hundreds of concurrent matches without stalling or hanging the primary server thread.

\subsection{Isolation Integrity and Sandbox Resilience Results}
Security sandboxing tests generated a 100\% success rate in blocking host-level compromise. The architecture was subjected to rigorous stress tests involving adversarial algorithms explicitly attempting to breach their partitioned environment via "Fork Bomb" scripts, infinite recursion loops, and aggressive memory consumption patterns. In every instance, the Linux Kernel \texttt{cgroups} consistently intervened. The Docker daemon gracefully terminated the malicious container the exact millisecond the 512MB RAM threshold or the 1.0-second CPU Turn limit was exceeded. Because these terminations were managed at the kernel level, the primary Flask server, the Redis queues, and other concurrently running tenant matches suffered absolutely zero downtime or noticeable performance degradation.

\subsection{Scalability and Resource Utilization Metrics}
Beyond security, the system's resource footprint was heavily monitored to evaluate its deployment scope. Because the platform relies heavily on localized processing, minimizing the CPU and Memory overhead per match was essential. The metrics proved that by stripping the Docker images down to mere Megabytes using Alpine Linux base images, the server could host significantly more concurrent execution matches than initially anticipated. During peak simulated loads, the CPU successfully balanced over 30 simultaneous background Python workers, delegating the IO-bound pipe reading effectively while saving processing cycles exclusively for the discrete bot computations.

\section{ARCHITECTURAL INNOVATIONS AND DISCUSSION}
The true distinction of the EMERGENT platform—and the core subjects of this discussion—lies in its advanced backend orchestration mechanics. These architectural novelties effectively elevate the platform from a simple competitive coding application to a robust, highly scalable, and structurally impenetrable execution framework.

\subsection{Intelligent Matchmaking via Max-Information-Gain}
A critical innovation in the system design is the implementation of the "Max-Information-Gain" scheduler policy. Standard round-robin or completely random matchmaking algorithms waste incredible amounts of server compute power pairing Grandmaster-level algorithms against basic, randomized beginner scripts. The results of such heavily skewed matches are statistically meaningless, as the outcome is effectively guaranteed before execution begins. EMERGENT introduces an intelligent queue sweeper that dynamically calculates the standard deviation of all active participants, explicitly pairing bots whose predicted win probability sits nearest to 50\%. This ensures that every computational cycle and container spun up by the server yields the highest possible mathematical refinement of the Glicko-2 ratings, drastically optimizing CPU usage and accelerating the tournament's leaderboard convergence to absolute truth.

\subsection{Asynchronous Backend Orchestration and Queue Management}
The complete decoupling of the web-facing REST handlers from the physical compilation engines is a massive architectural triumph. By dropping user code immediately into a memory-mapped Redis-backed Publish/Subscribe queue, the Flask server remains entirely unblocked and stateless. The Worker layer acts as an autonomous daemon, continuously polling for jobs, provisioning the necessary Inter-Process Communication (IPC) pipes, and resolving the matches entirely independently from the API process. This event-loop mechanism allows for seamless horizontal scaling—meaning tournament organizers could easily attach three, five, or ten worker machines to the same Redis instance without changing a single line of application code, ensuring a perfect separation of concerns.

\subsection{Zero-Trust Deterministic Sandboxing and Reproducibility}
Safely compiling and evaluating completely arbitrary user code is an inherently dangerous operation in software engineering. The platform's solution relies on aggressive, isolated containerization tied strictly to deterministic input processing. In the world of competitive algorithm design, reproducibility is paramount. The platform guarantees that two identical bot binaries playing on the exact same topological board seed will consistently yield a bit-for-bit identical \texttt{GameState.json} outcome without exception. This marriage of Absolute Determinism and Zero-Trust Isolation ensures that regardless of the chaotic nature of the submitted algorithms, competitive integrity remains completely uncompromised. Bugs, infinite loops, and timeout exceptions are perfectly reproducible and strictly confined to the responsible bot.

\subsection{Educational Impact vs Standard Static Analyzers}
The resulting platform completely disrupts the traditional evaluation paradigms heavily relied upon by modern educational institutions. While static analyzers (such as LeetCode tests) provide excellent foundational teaching metrics, they operate in completely sterile vacuums. Through the results observed in EMERGENT's live testing phases, developers who excelled at solving static mathematical arrays often struggled initially to handle the dynamic, state-driven chaos of an adversarial Tron match. The platform forces students to transition from writing rigid, sequential scripts to architecting highly resilient, asynchronous state machines. By penalizing algorithms that fail to gracefully handle exceptions or unexpected enemy strategies, EMERGENT acts as a far superior simulator for the actual demands of modern, real-world Cloud and Software Engineering production environments.

% ====================
% CHAPTER 7: CONCLUSION AND FUTURE SCOPE
% ====================
\chapter{CONCLUSION AND FUTURE SCOPE}

\section{CONCLUSION}
The EMERGENT platform successfully delivers a robust, autonomous LAN coding environment that elevates traditional algorithmic assessment into a dynamic, adversarial sandbox. By integrating a React.js/Monaco frontend with a localized Flask API, the system eliminates overhead while providing a frictionless developer experience. At its core, the Docker-backed execution engine ensures absolute zero-trust security—effectively isolating host systems from malicious or runaway code by enforcing strict compute limits (`cgroups`, 512MB RAM, 1.0s CPU Turn time).

Furthermore, the implementation of the Glicko-2 rating system combined with a "Max-Information-Gain" scheduler optimizes computational resources by pairing bots of similar skill levels, guaranteeing mathematically valid leaderboard convergence. The generated \texttt{GameState.json} replays empower users with visual, chronological debugging capabilities, fully resolving the "Black Box" nature of standard competitive platforms. Ultimately, the MVP meets all predefined objectives, providing a scalable and highly deterministic framework for automated programming evaluation.

\section{FUTURE SCOPE}
While the current minimum viable product (MVP) architecture is exceptionally functional for localized and LAN-based tournament configurations, the underlying framework was explicitly designed with modularity in mind. This foresight ensures that the platform is not severely bounded by its current implementation and naturally supports comprehensive extensions. The following domains have been identified as primary avenues for future enhancement:

\subsection{Expansion of the Game Plugin Registry}
The core logic engine currently operates via a robust but generalized Umpire protocol, cleanly demonstrated by the deterministic two-player `Tron` game model. Future iterations of the EMERGENT platform will expand this registry to include a diverse, multi-tiered array of topological and strategic environments. This includes integrating imperfect information games—such as variants of 'Battleship' with fog-of-war mechanics—which would force developer bots to rely heavily on probability matrices and state-estimation heuristics rather than absolute deterministic pathfinding. Additionally, the arena structure could be upgraded to handle multi-agent ($N\text{v}N$) entity environments where participants must program multiple sub-bots that concurrently coordinate and communicate via simulated local networking channels.

\subsection{Advanced Machine Learning Integration and Analytics}
The platform currently generates vast quantities of highly structured data in the form of raw textual execution outputs and precise spatial coordinate choices logged within the \texttt{GameState.json} artifacts. By integrating a dedicated data analytics pipeline, these historical database replays could be aggregated and leveraged to train generalized Machine Learning models. Implementing Deep Reinforcement Learning algorithms over thousands of matches could effectively act as a training ground for an infinitely scalable "Benchmark Agent"—a bot organically derived from aggregating and refining the absolute best algorithmic strategies devised by human participants.

\subsection{Team-Based Collaborative Workspaces}
To foster a collaborative, team-oriented development environment, the frontend implementation of the Monaco IDE can be significantly upgraded. By integrating Conflict-Free Replicated Data Types (CRDTs) or established Operational Transformation (OT) models, the editor could support real-time, synchronized multi-user programming environments. This would allow multiple students to connect their client-side browsers to the exact same bot profile simultaneously. Consequently, teammates could fluidly write completely separate sections of an algorithm—such as delegating target prioritization to one developer block and pathfinding to another—concurrently in real-time, greatly expediting team tournaments.

% ====================
% BIBLIOGRAPHY
% ====================
\cleardoublepage
\addcontentsline{toc}{chapter}{BIBLIOGRAPHY}
\begin{thebibliography}{99}

\bibitem{spa_react}
C. Mikkonen and T. Taivalsaari, ``Web apps vs. native apps: what is the future?,'' \emph{7th International Conference on Mobile Web Information Systems (MobiWIS 2010)}, pp. 211-216, 2010.

\bibitem{docker_paper}
D. Merkel, ``Docker: lightweight Linux containers for consistent development and deployment,'' \emph{Linux Journal}, vol. 2014, no. 239, p. 2, 2014.

\bibitem{docker_docs}
Docker Inc., \emph{Docker Official Documentation}. [Online]. Available: \url{https://docs.docker.com/}

\bibitem{redis_db}
J. Carlson, \emph{Redis in Action}. Shelter Island, NY: Manning Publications, 2013.

\bibitem{glicko_paper}
M. E. Glickman, ``Parameter estimation in large dynamic paired comparison experiments,'' \emph{Journal of the Royal Statistical Society: Series C (Applied Statistics)}, vol. 48, no. 3, pp. 377-394, 1999.

\bibitem{react_docs}
Meta Platforms, Inc., \emph{React Documentation}. [Online]. Available: \url{https://react.dev/}

\bibitem{monaco_editor}
Microsoft, \emph{Monaco Editor Repository}. [Online]. Available: \url{https://microsoft.github.io/monaco-editor/}

\bibitem{auto_grader}
P. Ihantola, T. Ahoniemi, V. Karavirta, and O. Seppälä, ``Review of recent systems for automatic assessment of programming assignments,'' \emph{Proceedings of the 10th Koli Calling International Conference on Computing Education Research}, pp. 86-93, 2010.

\bibitem{flask_docs}
Pallets, \emph{Flask Official Documentation}. [Online]. Available: \url{https://flask.palletsprojects.com/}

\bibitem{redis_docs}
Redis Ltd., \emph{Redis Official Documentation}. [Online]. Available: \url{https://redis.io/docs/}

\end{thebibliography}

\end{document}
